\section{Mathematik}

% ─────────────────────────────────────────────────────────────────────────────
\subsection{Abbildungen und deren Eigenschaften}

\definition{Eine \textbf{Abbildung} (oder Funktion) $f: X \to Y$ ordnet jedem Element $x$ aus dem
\textbf{Definitionsbereich} $X$ genau ein Element $f(x)$ aus dem \textbf{Wertebereich} $Y$ zu.}

\textbf{Wichtige Teilmengen:}
\begin{itemize}[leftmargin=1.5em]
    \item \textbf{Bildmenge} $f(X) = \{ f(x) \mid x \in X \} \subseteq Y$: die Menge aller tatsächlich angenommenen Werte.
    \item \textbf{Kern} $\ker(f) = \{ x \in X \mid f(x) = 0 \}$: Urbildmenge der Null (relevant bei linearen Abbildungen).
\end{itemize}

\textbf{Eigenschaften:}
\begin{itemize}[leftmargin=1.5em]
    \item \textbf{Injektiv} (eineindeutig): $f(x_1) = f(x_2) \Rightarrow x_1 = x_2$. Jedes Element der Bildmenge hat genau ein Urbild. Für lineare Abbildungen äquivalent zu $\ker(f) = \{0\}$.
    \item \textbf{Surjektiv} (auf): $f(X) = Y$, d.\,h.\ jedes $y \in Y$ wird angenommen.
    \item \textbf{Bijektiv}: injektiv \emph{und} surjektiv - Voraussetzung für die Existenz der \textbf{Umkehrabbildung} $f^{-1}: Y \to X$.
    \item \textbf{Monoton wachsend}: $x_1 < x_2 \Rightarrow f(x_1) \leq f(x_2)$; \textbf{streng monoton wachsend} mit $<$ statt $\leq$ (impliziert Injektivität auf $\mathbb{R}$).
    \item \textbf{Fixpunkte}: $x^*$ mit $f(x^*) = x^*$; relevant z.\,B.\ bei Iterationsverfahren und dynamischen Systemen.
\end{itemize}

\textbf{Umkehrbarkeit:} Eine Abbildung ist genau dann umkehrbar, wenn sie bijektiv ist.
Bei linearen Abbildungen $f(\bm{x}) = A\bm{x}$ ist dies äquivalent zu $\det(A) \neq 0$ (reguläre Matrix).

\beispiel{$f: \mathbb{R} \to \mathbb{R},\; f(x) = x^2$ ist weder injektiv (da $f(1)=f(-1)$) noch surjektiv (negative Werte werden nicht angenommen) und damit nicht umkehrbar. Eingeschränkt auf $[0,\infty)$ ist $f$ bijektiv und $f^{-1}(y) = \sqrt{y}$.}

\newpage
% ─────────────────────────────────────────────────────────────────────────────
\subsection{Vektor- und Matrixnormen}

\definition{Eine \textbf{Norm} $\|\cdot\|$ auf einem Vektorraum ist eine Abbildung in $\mathbb{R}_{\geq 0}$, die folgende vier Axiome erfüllt:
\begin{enumerate}[leftmargin=2em]
    \item $\|\bm{v}\| \geq 0$ (Nichtnegativität)
    \item $\|\bm{v}\| = 0 \Leftrightarrow \bm{v} = \bm{0}$ (Definitheit)
    \item $\|\alpha\bm{v}\| = |\alpha|\cdot\|\bm{v}\|$ (Homogenität)
    \item $\|\bm{u}+\bm{v}\| \leq \|\bm{u}\|+\|\bm{v}\|$ (Dreiecksungleichung)
\end{enumerate}}

\textbf{Vektornormen ($p$-Normen):}
\formel{\|\bm{v}\|_p = \left(\sum_{i=1}^n |v_i|^p\right)^{1/p}}
Wichtige Spezialfälle: $\|\bm{v}\|_1 = \sum|v_i|$ (Betragssumme), $\|\bm{v}\|_2 = \sqrt{\sum v_i^2}$ (euklidische Norm), $\|\bm{v}\|_\infty = \max_i|v_i|$ (Maximumsnorm).

\textbf{Matrixnormen:}
\begin{itemize}[leftmargin=1.5em]
    \item \textbf{Höldersche-$p$-Norm} (elementweise): $\|A\|_{Hp} = \left(\sum_{i,j}|a_{ij}|^p\right)^{1/p}$
    \item \textbf{Frobenius-Norm} ($= H_2$-Norm $=$ Schatten-2-Norm): $\|A\|_F = \sqrt{\sum_{i,j} a_{ij}^2} = \sqrt{\mathrm{Spur}(A^\top A)}$
    \item \textbf{Spektralnorm} ($= $ Schatten-$\infty$-Norm): $\|A\|_S = \sqrt{\rho(A^\top A)} = \sigma_{\max}(A)$ (größter Singulärwert)
\end{itemize}

\textbf{Submultiplikativität:} Für induzierte Normen gilt $\|AB\| \leq \|A\|\cdot\|B\|$ - wichtig für Stabilitätsaussagen.

\textbf{Zusammenhang mit Spektralradius und Dynamik:} Für den Spektralradius $\rho(A) = \max_i |\lambda_i|$ gilt:
\formel{\rho(A) \leq \|A\|_S \leq \|A\|_F}
Ein Pulsmodell ist \emph{kontrahierend} (kurzfristige Dämpfung), wenn $\|A\|_S < 1$; \emph{konvergent} (langfristige Dämpfung), wenn $\rho(A) < 1$. Jedes kontrahierende Modell ist konvergent, aber nicht umgekehrt.

\textbf{Kondition einer Matrix:}
\formel{\kappa_S(A) = \|A\|_S \cdot \|A^{-1}\|_S = \frac{\sigma_{\max}(A)}{\sigma_{\min}(A)}}
Je größer die Kondition, desto stärker werden Fehler bei numerischen Verfahren verstärkt. Optimale Kondition = 1 (Einheitsmatrix / orthogonale Matrizen).

\newpage
% ─────────────────────────────────────────────────────────────────────────────
\subsection{Eigen- und Singulärwerte sowie Eigen- und Singulärvektoren und deren geometrische Interpretation}

\definition{Ein \textbf{Eigenvektor} $\bm{e}_i \neq \bm{0}$ einer Matrix $A$ und der zugehörige \textbf{Eigenwert} $\lambda_i \in \mathbb{C}$ erfüllen:
\formel{A \cdot \bm{e}_i = \lambda_i \cdot \bm{e}_i}
Der Eigenvektor ändert durch die Multiplikation mit $A$ seine \emph{Richtung nicht} (bei negativem $\lambda_i$ wird die Orientierung umgekehrt); er wird lediglich um den Faktor $\lambda_i$ gestreckt oder gestaucht.}

\textbf{Singulärwerte und -vektoren:} Die \textbf{Singulärwerte} $\sigma_i(A) = \sqrt{\lambda_i(A^\top A)} \geq 0$ sind die Wurzeln der Eigenwerte von $A^\top A$. Die zugehörigen \textbf{Singulärvektoren} stehen paarweise \emph{orthogonal} aufeinander (im Gegensatz zu Eigenvektoren).

\textbf{Geometrische Interpretation (2D):}
Eine $2 \times 2$-Matrix $A$ transformiert alle Vektoren des Einheitskreises in eine \textbf{Ellipse}.
\begin{itemize}[leftmargin=1.5em]
    \item Die \textbf{Eigenvektoren} sind jene Richtungen, die unter $A$ \emph{unverändert} bleiben - sie zeigen auf die Punkte des Bildes, die noch auf denselben Strahlen liegen wie im Urbild.
    \item Die \textbf{Singulärvektoren} beschreiben die \textbf{Halbachsen der resultierenden Ellipse}; die Singulärwerte geben deren Längen an.
\end{itemize}

\textbf{Zusammenhang Determinante und Singulärwerte:}
\formel{|\det(A)| = \prod_i \sigma_i(A)}
Der Absolutbetrag der Determinante ist gleich dem Produkt aller Singulärwerte und beschreibt damit das Verhältnis des Volumens des transformierten Ellipsoids zur Einheitskugel.

\textbf{Vorzeichen der Determinante:} $\det(A) > 0$ bedeutet, dass die Orientierung erhalten bleibt; $\det(A) < 0$ bedeutet, dass die Transformation eine Spiegelung enthält.

\newpage
% ─────────────────────────────────────────────────────────────────────────────
\subsection{Eigenwertproblem und Diagonalisierung von quadratischen Matrizen}

\textbf{Charakteristisches Polynom:} Die Eigenwerte $\lambda_i$ einer $n \times n$-Matrix $A$ sind die Nullstellen des charakteristischen Polynoms:
\formel{P_n(\lambda) = \det(A - \lambda E) = 0}
Das Polynom hat Grad $n$ und liefert (über $\mathbb{C}$) genau $n$ Eigenwerte (mit Vielfachheit).

\textbf{Eigenschaften:}
\begin{itemize}[leftmargin=1.5em]
    \item $\mathrm{Spur}(A) = \sum_i \lambda_i$ (Summe der Hauptdiagonalelemente = Summe der Eigenwerte)
    \item $\det(A) = \prod_i \lambda_i$ (Determinante = Produkt der Eigenwerte)
    \item Ist ein Eigenwert $\lambda_i = 0$, so ist $A$ \emph{singulär} ($\det(A) = 0$, nicht invertierbar).
\end{itemize}

\textbf{Diagonalisierung:} Eine $n \times n$-Matrix $A$ ist \textbf{diagonalisierbar}, wenn sie $n$ linear unabhängige Eigenvektoren besitzt. Dann gilt:
\formel{A = V \Lambda V^{-1}, \quad \Lambda = \mathrm{diag}(\lambda_1, \ldots, \lambda_n)}
wobei die Spalten von $V$ die Eigenvektoren sind. Matrizen mit $n$ verschiedenen Eigenwerten sind stets diagonalisierbar.

\textbf{Defekte Matrizen:} Eine Matrix heißt \textbf{defekt}, wenn die geometrische Vielfachheit (Dimension des Eigenraums) für einen Eigenwert kleiner als dessen algebraische Vielfachheit ist. Defekte Matrizen sind \emph{nicht diagonalisierbar}.

\beispiel{Für eine $2 \times 2$-Matrix $A = \begin{pmatrix} a & b \\ c & d \end{pmatrix}$ lautet das charakteristische Polynom:
$\lambda^2 - (a+d)\lambda + (ad-bc) = 0$, lösbar mit der Lösungsformel $\lambda_{1,2} = \frac{(a+d) \pm \sqrt{(a+d)^2 - 4(ad-bc)}}{2}$.}

\newpage
% ─────────────────────────────────────────────────────────────────────────────
\subsection{Fuzzy-Logik, Fuzzy-Sets und unscharfe Zahlen}

\textbf{Motivation - Grenzen der klassischen Logik:} Die zweiwertige Logik (wahr/falsch) stößt bei der Modellierung gradueller Übergänge an ihre Grenzen (Sorites-Paradoxon: „Wann ist ein Sandhaufen kein Sandhaufen mehr?"). Zudem führt die Implikation $f \to *$ in der klassischen Logik immer zu \emph{wahr} (\emph{ex falso quodlibet}), was in vielen Kontexten kontraintuitiv ist.

\definition{In der \textbf{Fuzzy-Logik} (L.\,A.\ Zadeh, 1965) werden die Wahrheitswerte von $\{0,1\}$ auf das Kontinuum $[0,1]$ verallgemeinert. Logische Verknüpfungen werden z.\,B.\ definiert als:
\formel{\llbracket \lnot A \rrbracket = 1 - \llbracket A \rrbracket, \quad \llbracket A \land B \rrbracket = \min(\llbracket A \rrbracket, \llbracket B \rrbracket), \quad \llbracket A \lor B \rrbracket = \max(\llbracket A \rrbracket, \llbracket B \rrbracket)}
Weitere Varianten sind \textbf{t-Normen} (für $\land$) und \textbf{t-Conormen} (für $\lor$), z.\,B.\ das Produkt $T_{\mathrm{prod}}(x,y) = x \cdot y$.}

\definition{Eine \textbf{Fuzzy-Menge} $M \subseteq X$ ist definiert durch eine \textbf{Zugehörigkeitsfunktion} $\mu_M: X \to [0,1]$. Der Wert $\mu_M(x)$ gibt den \emph{graduellen Zugehörigkeitsgrad} von $x$ zu $M$ an. Scharfe (klassische) Mengen sind der Spezialfall mit $\mu_M(x) \in \{0,1\}$.}

\textbf{Mengenoperationen} werden über die Fuzzy-Logik definiert:
$\mu_{M^c}(x) = 1 - \mu_M(x)$,\quad $\mu_{A \cap B}(x) = \min(\mu_A(x), \mu_B(x))$,\quad $\mu_{A \cup B}(x) = \max(\mu_A(x), \mu_B(x))$.

\textbf{$\alpha$-Schnitte:} $[\mu]_\alpha = \{x \in X \mid \mu(x) \geq \alpha\}$ - jene Teilmenge, deren Zugehörigkeitsgrad mindestens $\alpha$ beträgt. Ermöglicht eine algorithmisch handhabbare Darstellung von Fuzzy-Mengen.

\textbf{Unscharfe Zahlen und Fuzzy-Arithmetik:} Fuzzy-Mengen können \emph{unscharfe Zahlen} (z.\,B.\ „ungefähr 2") darstellen. Rechenoperationen auf unscharfen Zahlen sind möglich, führen jedoch i.\,A.\ zu \emph{wachsender Unschärfe} und sind nicht exakt umkehrbar (z.\,B.\ $({\sim}2) + ({\sim}2) - ({\sim}2) \neq {\sim}2$).

\newpage
% ─────────────────────────────────────────────────────────────────────────────
\subsection{Modulare Arithmetik und Integer-Darstellung am Computer}

\textbf{Zahlendarstellung:} Jede ganze Zahl lässt sich in einer beliebigen Basis $b \geq 2$ darstellen. Wichtige Basen: Dezimal ($b=10$), Binär ($b=2$, physikalisch realisierbar durch EIN/AUS), Oktal ($b=8$), Hexadezimal ($b=16$, Umrechnung in 4-Bit-Blöcken). Umrechnung Dezimal $\to$ Binär: sukzessive Division durch 2, Reste rückwärts lesen.

\textbf{Ganzzahl-Darstellung am Computer:} Mit $n$ Bits lassen sich $2^n$ verschiedene Werte darstellen. Bei vorzeichenbehafteten Integern (Zweierkomplement) ergibt sich z.\,B.\ für 16 Bit der Bereich $[-32768, 32767]$. \textbf{Overflow} tritt auf, wenn ein Ergebnis diesen Bereich verlässt (z.\,B.\ $17000 + 18000 = 35000 > 32767$ ergibt fälschlicherweise $-30536$).

\textbf{Teilbarkeit und Primzahlen:} $m \mid n$ (``$m$ teilt $n$''), wenn $n = k \cdot m$ für ein $k \in \mathbb{Z}$. \textbf{Primzahlen} $p \in \{2,3,5,7,11,\ldots\}$ haben nur die trivialen Teiler $1$ und $p$. Jede natürliche Zahl $n \geq 2$ besitzt eine eindeutige \textbf{Primfaktorzerlegung}.

\definition{\textbf{Modulare Arithmetik:} $a \equiv b \pmod{n}$ (kongruent modulo $n$), wenn $n \mid (a-b)$, d.\,h.\ $a$ und $b$ haben denselben Rest bei Division durch $n$. Die Menge der Restklassen $\mathbb{Z}/n\mathbb{Z} = \{0,1,\ldots,n-1\}$ bildet den \textbf{Restklassenring}. Addition, Subtraktion und Multiplikation sind mit der Kongruenz verträglich.}

\textbf{Euklidischer Algorithmus:} Berechnet $\gcd(a,b)$ effizient durch sukzessive Modulo-Reduktion:
$\gcd(a,b) = \gcd(a \bmod b,\, b)$, bis der Rest 0 wird. Der \textbf{erweiterte euklidische Algorithmus} liefert zusätzlich $\alpha, \beta \in \mathbb{Z}$ mit $\alpha a + \beta b = \gcd(a,b)$ (Bézout-Koeffizienten).

\textbf{Multiplikatives Inverses:} Das multiplikative Inverse von $a$ modulo $n$ ist jenes $b \in \mathbb{Z}/n\mathbb{Z}$ mit $a \cdot b \equiv 1 \pmod{n}$. Es existiert genau dann, wenn $\gcd(a,n) = 1$. Bei Primzahl $p$ als Modul existiert das multiplikative Inverse für alle $a \neq 0$ - $\mathbb{Z}/p\mathbb{Z}$ bildet damit einen \textbf{Körper}.

\textbf{Hamming-Distanz:} Zwischen zwei Bitfolgen gleicher Länge $b_1, b_2$ ist $d_H(b_1,b_2)$ die Anzahl der Stellen, an denen sie sich unterscheiden. Sie erfüllt alle Metraxiome und ist gleich dem \textbf{Hamming-Gewicht} $w(b_1 \oplus b_2)$ (Anzahl der Einsen in der XOR-Summe). Grundlage für fehlererkennende und fehlerkorrigierende Codes (z.\,B.\ Paritätsbit in ASCII, Hamming-Code).

\newpage


\subsection{Darstellung reeller Zahlen am Computer}
 
Am Computer steht zur Darstellung einer einzelnen Zahl nur eine
\textbf{begrenzte Anzahl von Bits} zur Verfügung.
Für reelle Zahlen gibt es zwei grundlegende Ansätze:
 
\medskip
\textbf{Festkommadarstellung}
\begin{itemize}
  \item Eine feste Anzahl von Bits vor und nach dem Komma.
  \item Gleichmäßige Verteilung der Zahlen im Wertebereich
        $\Rightarrow$ absoluter Fehler überall gleich groß.
  \item \emph{Nachteile:} geringer Wertebereich; bei kleinen Zahlen ist
        der absolute Fehler relativ groß (z.\,B. $\Delta x = 0{,}01$
        bei Werten im Bereich $0{,}1$--$0{,}9$ ungünstig).
  \item Heute selten verwendet.
\end{itemize}
 
\medskip
\textbf{Gleitkommadarstellung (Floating Point)}
 
\begin{definitionbox}{IEEE-754 Gleitkommazahl}
Eine Gleitkommazahl hat die Form
\[
  x = \pm\; 1{,}\underbrace{b_1 b_2 \cdots b_m}_{\text{Mantisse}}
      \;\times\; 2^{\,\pm e},
\]
bestehend aus \textbf{Vorzeichen}, \textbf{Mantisse} und \textbf{Exponent}.
Das führende \texttt{1.} muss nicht gespeichert werden (\emph{Hidden Bit}).
 
\medskip
\begin{tabular}{lcccc}
  \toprule
  Format & Bits gesamt & Vorzeichen & Exponent & Mantisse \\
  \midrule
  Single Precision & 32 & 1 & 8 & 23 \\
  Double Precision & 64 & 1 & 11 & 52 \\
  \bottomrule
\end{tabular}
\end{definitionbox}
 
\begin{itemize}
  \item \textbf{Sehr großer Wertebereich} (von sehr kleinen bis sehr großen Zahlen).
  \item \textbf{Gleich viele Zahlen pro Abschnitt} (auf logarithmischer Skala):
        nahe der~$0$ liegen die darstellbaren Zahlen dichter als bei großen Werten.
  \item Innerhalb des Wertebereichs nur \emph{näherungsweise} Darstellung möglich
        $\Rightarrow$ Rundungsfehler.
\end{itemize}
 
\subsection{Numerische Fehler}
 
Rechnen mit begrenzter Genauigkeit führt im Allgemeinen zu Fehlern.
Das Ziel ist stets, jede Zahl durch die nächste \textbf{Maschinenzahl}
darzustellen -- doch auch die Summe bester Darstellungen ist nicht
mehr unbedingt die beste Darstellung des Ergebnisses.
 
\begin{examplebox}{Motivationsbeispiel}
$\frac{1}{3}+\frac{1}{3}+\frac{1}{3}=1$, aber in endlicher Darstellung:
$0{,}333 + 0{,}333 + 0{,}333 = 0{,}999 \neq 1{,}000$.
 
Ebenso: $0{,}1_{10} = 0{,}\overline{0011}_2$ -- Dezimalzahlen sind im
Binärsystem meist nur mit Fehler darstellbar.
\end{examplebox}
 
\medskip
\textbf{Unangenehme Fehlerarten:}
 
\begin{definitionbox}{Auslöschung}
Subtraktion zweier fast gleich großer Zahlen $\Rightarrow$ Ergebnis wird
sehr ungenau.\\[4pt]
\textbf{Beispiel:}
\[
  \underbrace{1{,}234567}_{\text{7 sig. Stellen}}
  - \underbrace{1{,}234565}_{\text{7 sig. Stellen}}
  = \underbrace{0{,}000002}_{\text{nur noch 1 sig. Stelle}}
  = 2{,}000000 \times 10^{-6}
\]
\textbf{Regel:} Wenn möglich, keine kleine Zahl aus der Differenz zweier
großer Zahlen berechnen!
\end{definitionbox}
 
\begin{definitionbox}{Absorption}
Addition einer Zahl zu einer betragsmäßig viel größeren kann keine
Veränderung bewirken.\\[4pt]
\textbf{Beispiel} (7 Stellen Genauigkeit):
\[
  1\,245\,327 + 0{,}215231 \;\approx\; 1\,245\,327
  \quad\Rightarrow\quad \text{keine Änderung!}
\]
\textbf{Gefahr} bei wiederholter Addition kleiner Inkremente
$\Rightarrow$ gesamte Addition hat keine Wirkung.
\end{definitionbox}
 
\begin{definitionbox}{Unterlauf (Underflow)}
Multiplikation betragsmäßig kleiner Zahlen
$\Rightarrow$ Ergebnis wird am besten durch $0$ dargestellt,
obwohl das mathematisch falsch ist.\\[4pt]
\textbf{Beispiel:}
\[
  10^{-70} \cdot 10^{-70} \cdot 10^{70} \cdot 10^{70}
  \;\to\; 0 \cdot 10^{70} \cdot 10^{70} = 0
  \quad\text{(falsch, da } \approx 1\text{)}
\]
\end{definitionbox}
 
% ═══════════════════════════════════════════════════════════════════════════
\section{Frage 8 · Kondition von Problemen}
% ═══════════════════════════════════════════════════════════════════════════
 
\subsection{Grundlegende Konzepte}
 
\begin{definitionbox}{Kondition eines Problems}
Die \textbf{Kondition} beschreibt die prinzipielle \emph{Gutartigkeit}
(oder \emph{Bösartigkeit}) eines Problems selbst -- unabhängig vom
gewählten Algorithmus.
 
\medskip
\textbf{Kernfrage:} Wie empfindlich ist das Ergebnis gegenüber
kleinen Störungen in den Eingangsdaten (Messunsicherheiten,
Rundungsfehler)?
 
\medskip
\begin{tabular}{ll}
  \textbf{gut konditioniert:} & kleine Störung $\Rightarrow$ geringe Ergebnisänderung \\
  \textbf{schlecht konditioniert:} & kleine Störung $\Rightarrow$ große Ergebnisänderung \\
\end{tabular}
 
\medskip
Die Kondition ist eine \emph{Untergrenze des zu erwartenden Fehlers} --
mit einem schlecht gewählten Algorithmus kann es noch schlimmer kommen.
\end{definitionbox}
 
\subsection{Absoluter und relativer Fehler}
 
Betrachte eine Abbildung $f: X \to Y$, $x \mapsto y = f(x)$, wobei in
$X$ und $Y$ jeweils eine Norm definiert ist.
 
Liegt eine Störung $\delta x$ vor (d.\,h. $x \to x + \delta x$),
so wird auch der Ausgabewert gestört: $y \to y + \delta y$.
 
\begin{definitionbox}{Absoluter und relativer Fehler}
\[
  \|\delta x\|_X, \;\|\delta y\|_Y
  \quad\ldots\quad \text{absolute Fehler}
\]
\[
  \delta_x = \frac{\|\delta x\|_X}{\|x\|_X}, \quad
  \delta_y = \frac{\|\delta y\|_Y}{\|y\|_Y}
  \quad\ldots\quad \text{relative Fehler}
\]
Relative Fehler sind oft aussagekräftiger, da sie skaliert sind
(sie gehen erst auf Verhältnisskala).
\end{definitionbox}
 
\subsection{Konditionszahl}
 
\begin{definitionbox}{Relative Konditionszahl $\krel$}
Die relative Konditionszahl ist die maximale Fehlerverstärkung im
Worst Case für sehr kleine Störungen:
\[
  \krel = \limsup_{\delta_x \to 0} \frac{\delta_y}{\delta_x}
\]
\textbf{Für differenzierbare Abbildungen (1-dimensional):}
\[
  \krel(x) = \left|\frac{f'(x) \cdot x}{f(x)}\right|
\]
Diese Zahl ist \emph{ortsabhängig}.
 
\textbf{Mehrdimensional} (mit Maximumsnorm, $f: \mathbb{R}^n \to \mathbb{R}$):
\[
  \krel(\vx) = \max_{i=1,\ldots,n}
  \left|\frac{\partial f}{\partial x_i}(\vx)\cdot\frac{x_i}{f(\vx)}\right|
\]
 
\textbf{Allgemein:}
\[
  \krel(x) = \frac{\|\nabla h(x)\|}{|h(x)|} \cdot \|x\|
\]
\end{definitionbox}
 
\subsection{Wichtige Spezialfälle}
 
\textbf{Kondition der Addition} ($f_+(x_1, x_2) = x_1 + x_2$):
\[
  \frac{\partial f_+}{\partial x_i} = 1
  \;\Rightarrow\;
  \krel = \max\!\left\{\frac{|x_1|}{|x_1+x_2|},\, \frac{|x_2|}{|x_1+x_2|}\right\}
\]
\begin{itemize}
  \item Gleiche Vorzeichen: $\krel \le 1$ $\Rightarrow$ keine Fehlerverstärkung.
  \item Ungleiche Vorzeichen: $\krel$ kann groß werden
        $\Rightarrow$ \textbf{Fehlerverstärkung!}
  \item Extremfall $x_1 \approx -x_2$: $\krel \gg 1$
        $\Rightarrow$ dies ist der Effekt der \textbf{Auslöschung}.
\end{itemize}
 
\textbf{Kondition der Multiplikation} ($f_*(x_1, x_2) = x_1 x_2$):
\[
  \frac{\partial f_*}{\partial x_1} = x_2, \quad
  \frac{\partial f_*}{\partial x_2} = x_1
  \;\Rightarrow\;
  \krel = 1
\]
$\Rightarrow$ Multiplikation ist \textbf{gut konditioniert}. Einzige
Gefahr ist Über- bzw.\ Unterlauf (Verlassen des darstellbaren Bereichs).
 
\medskip
\textbf{Kondition linearer Gleichungssysteme} ($y = \mA x$):
\begin{itemize}
  \item Die Konditionszahl einer Matrix ergibt sich als
        Verhältnis des größten zum kleinsten Singulärwert:
        \[
          \mathrm{cond}(\mA) = \frac{\sigma_{\max}}{\sigma_{\min}}
          = \|\mA\|_S \cdot \|\mA^{-1}\|_S \;\ge 1
        \]
  \item Je größer $\mathrm{cond}(\mA)$, desto \emph{schlechter} konditioniert.
  \item \emph{Anschaulich:} Fast parallele Geraden (fast singuläre Matrix)
        $\Rightarrow$ kleine Störung der Koeffizienten verschiebt den
        Schnittpunkt dramatisch.
  \item Die Determinante allein ist kein guter Indikator für die Kondition
        (sie kann durch Skalierung der Gleichungen beliebig verändert werden).
\end{itemize}
 
\begin{warnbox}{Achtung: Kondition vs.\ Stabilität}
Die Kondition ist eine Eigenschaft des \textbf{Problems}, nicht des
Algorithmus. Auch ein perfekter Algorithmus kann ein schlecht
konditioniertes Problem nicht besser lösen. Stabilität beschreibt
dagegen, ob ein Algorithmus den durch die Kondition vorgegebenen
Fehler nicht \emph{zusätzlich} vergrößert.
\end{warnbox}
 
% ═══════════════════════════════════════════════════════════════════════════
\section{Frage 9 · Minimumsuche: Gradientenabstieg, Newton-Verfahren, Impulsmethode}
% ═══════════════════════════════════════════════════════════════════════════
 
\subsection{Allgemeines Abstiegsverfahren}
 
Optimierungsaufgaben werden meist als Minimierungsprobleme formuliert.
Da eine direkte Lösung von $\nabla f = 0$ oft nicht möglich ist,
verwendet man iterative \textbf{Abstiegsverfahren}:
 
\begin{definitionbox}{Allgemeines Abstiegsschema}
\begin{enumerate}
  \item Wähle Startpunkt $\vx_0$.
  \item Wiederhole für $k = 0, 1, 2, \ldots$:
        \begin{enumerate}
          \item Ermittle Abstiegsrichtung $\vd_k$.
          \item Wähle Schrittweite $\lambda_k > 0$.
          \item Setze $\vx_{k+1} = \vx_k + \lambda_k \vd_k$.
        \end{enumerate}
  \item Abbruch bei Konvergenz oder maximaler Iterationsanzahl.
\end{enumerate}
Es muss gelten: $f(\vx_{k+1}) < f(\vx_k)$.
\end{definitionbox}
 
\subsection{Methode des steilsten Abstiegs (Gradientenabstieg)}
 
\begin{definitionbox}{Gradientenabstieg (Steepest Descent)}
Die Abstiegsrichtung ist der \textbf{negative Gradient}:
\[
  \vd_k = -\nabla f(\vx_k)
\]
Das mehrdimensionale Minimierungsproblem reduziert sich auf ein
\textbf{eindimensionales Linesearch-Problem}:
\[
  \lambda_k = \arg\min_{\lambda > 0} f\!\left(\vx_k + \lambda \vd_k\right)
\]
Für quadratische Funktionen $f(\vx) = \frac{1}{2}\vx^\top\mA\vx - \vx^\top\bm{b}$
gibt es eine analytische Formel für $\lambda_k$.
\end{definitionbox}
 
\textbf{Problem: Zick-Zack-Verhalten}
\begin{itemize}
  \item Aufeinanderfolgende Suchrichtungen stehen (idealerweise) orthogonal
        zueinander.
  \item Ziel ist es, jede Richtung ``maximal auszuschöpfen''.
  \item In der Praxis entsteht oft ein Zick-Zack-Kurs, besonders bei
        schlecht konditionierten Problemen (stark unterschiedliche Eigenwerte
        der Hesse-Matrix).
\end{itemize}
 
\textbf{Schrittweite: Armijo-Regel}
 
\begin{importantbox}{Armijo-Regel}
Wähle Parameter $q \in (0,1)$ und $\alpha \in (0, \tfrac{1}{2}]$.
Finde die größte Zahl $\lambda_k \in \{1, q, q^2, q^3, \ldots\}$ mit:
\[
  f(\vx_k + \lambda_k \vd_k) - f(\vx_k)
  \;\le\; \lambda_k\,\alpha\,\bigl(\nabla f(\vx_k)\bigr)^\top \vd_k
\]
Interpretation: Ausgehend von $\lambda^* = 1$ wird die Schrittweite
durch Multiplikation mit $q$ so lange reduziert, bis sie in den
``erlaubten Bereich'' fällt.
\end{importantbox}
 
\subsection{Konjugierte Gradienten (CG)}
 
\begin{definitionbox}{Conjugate Gradient Verfahren}
Behebt die Zick-Zack-Problematik des steilsten Abstiegs.
 
\textbf{Idee:} Wähle die neue Suchrichtung $\vd_{k+1}$
\textbf{$\mA$-orthogonal} (konjugiert) zur vorherigen:
\[
  \langle \vd_{k+1},\, \vd_k \rangle_{\mA}
  = \vd_{k+1}^\top\, \mA\, \vd_k = 0
\]
Für ein $n$-dimensionales Problem gibt es \emph{prinzipiell} nur $n$
unabhängige Richtungen $\Rightarrow$ nach \textbf{$n$ Schritten}
(bei quadratischen Problemen) ist die Lösung exakt gefunden.
 
\medskip
\textbf{Python:} \texttt{scipy.sparse.linalg.cg} (LGS) /
\texttt{scipy.optimize.minimize(method='CG')} (allg.)
\end{definitionbox}
 
\subsection{Newton-Abstieg}
 
Verfahren 2.\ Ordnung: nutzt Information aus der zweiten Ableitung
(Hesse-Matrix).
 
\begin{definitionbox}{Newton-Abstieg}
\textbf{Zentrale Formel} (für jeden Schritt $k = 0, 1, \ldots$):
\[
  \mH_f(\vx_k)\, \bm{s}_k = -\nabla f(\vx_k),
  \qquad \vx_{k+1} = \vx_k + \bm{s}_k
\]
Die Abstiegsrichtung $\bm{s}_k$ ergibt sich durch Lösung eines
\textbf{linearen Gleichungssystems} mit der Hesse-Matrix $\mH_f$.
 
\begin{itemize}
  \item Entspricht der Näherung der Zielfunktion durch ein Paraboloid
        und der Minimierung dieses Paraboloids.
  \item Konvergiert \emph{sehr rasch} (quadratisch), wenn der Startpunkt
        nahe genug am Minimum liegt.
  \item Kann gegen \textbf{Maxima oder Sattelpunkte} konvergieren
        (kein Wissen über Minimumsuche eingebaut!).
\end{itemize}
\end{definitionbox}
 
\textbf{Rechenaufwand pro Schritt:}
\begin{itemize}
  \item Gradient: $n$ Elemente.
  \item Hesse-Matrix: $\frac{n(n+1)}{2} \in \mathcal{O}(n^2)$ Elemente.
  \item LGS lösen: $\sim \frac{2}{3} n^3$ Operationen.
\end{itemize}
 
\textbf{Varianten zur Aufwandsreduktion:}
 
\begin{definitionbox}{Newton-Varianten}
\begin{itemize}
  \item \textbf{Frozen Newton:} Dieselbe Hesse-Matrix für mehrere
        aufeinanderfolgende Schritte verwenden. LU-Zerlegung erlaubt
        Reduktion von $\mathcal{O}(n^3)$ auf $\mathcal{O}(n^2)$
        für Folgeschritte.
  \item \textbf{Inexact Newton:} Lineares System nur näherungsweise
        lösen (z.\,B.\ erste Schritte des CG-Verfahrens).
  \item \textbf{BFGS} (\emph{Broyden--Fletcher--Goldfarb--Shanno}):
        Approximiere die Hesse-Matrix durch ein Update-Verfahren
        (Rang-2-Update) -- sehr populäres Quasi-Newton-Verfahren:
        \[
          \mB_{k+1} = \mB_k
          - \frac{(\mB_k\bm{s}_k)(\mB_k\bm{s}_k)^\top}{\bm{s}_k^\top(\mB_k\bm{s}_k)}
          + \frac{\bm{y}_k\bm{y}_k^\top}{\bm{s}_k^\top\bm{y}_k},
          \quad \bm{y}_k = \nabla f(\vx_{k+1}) - \nabla f(\vx_k)
        \]
\end{itemize}
\end{definitionbox}
 
\subsection{Line Search vs.\ Trust Region}
 
\begin{importantbox}{Zwei grundlegende Strategien}
\textbf{Line Search:}
\begin{enumerate}
  \item Bestimme Suchrichtung $\vd_k$ (z.\,B.\ steilster Abstieg).
  \item Suche Minimum (oder Näherung) entlang dieser Richtung.
  \item Wiederhole für den neuen Punkt.
\end{enumerate}
 
\textbf{Trust Region:}
\begin{enumerate}
  \item Nähere die Zielfunktion durch eine einfachere Modellfunktion
        (z.\,B.\ Polynom 2.\ Grades) an.
  \item Minimiere die Modellfunktion innerhalb eines Vertrauensbereichs.
  \item Bewerte die Qualität des gefundenen Punkts:
        \begin{itemize}
          \item gut $\Rightarrow$ Punkt annehmen, Vertrauensbereich vergrößern.
          \item schlecht $\Rightarrow$ Vertrauensbereich verkleinern, erneut suchen.
        \end{itemize}
\end{enumerate}
\end{importantbox}
 
\subsection{Abstiegsverfahren im Machine Learning}
 
\textbf{Besondere Herausforderungen} bei tiefen neuronalen Netzen:
\begin{itemize}
  \item Extrem viele Parameter (Gewichte $\mathbf{W}$).
  \item Sehr zerklüftete Verlustfunktion $L$ ohne schöne analytische Eigenschaften.
  \item Minimum i.\,A.\ nicht eindeutig (Ununterscheidbarkeitsproblem).
  \item Zielfunktion sicher nicht konvex.
\end{itemize}
 
\textbf{Lernraten statt exakter Liniensuche:}
\begin{definitionbox}{SGD mit Lernrate}
Da eine exakte Liniensuche ($\lambda_k = \arg\min f(\vx_k + \lambda\vd_k)$)
zu aufwändig ist, führt man eine \textbf{Lernrate} $\alpha_t$ ein:
\[
  \vx_{t+1} = \vx_t + \alpha_t\, \vd_t, \qquad
  \vd_t = -\nabla_{\mathbf{W}} L(\vx_t)
\]
Die Lernrate wird im Laufe des Trainings nach einer Vorschrift reduziert,
z.\,B.\ exponentiell $\alpha(t) = e^{-ct}$ oder hyperbolisch
$\alpha(t) = \frac{\alpha_0}{1+ct}$.
 
\medskip
\textbf{Stochastic Gradient Descent (SGD):} Um künstliche Zyklen zu
vermeiden, wird die Reihenfolge der Daten (bzw.\ die Zusammensetzung
der Mini-Batches) immer wieder zufällig verändert.
\end{definitionbox}
 
\subsection{Impulsmethode (Momentum)}
 
\begin{definitionbox}{Momentum-Methode}
Klassischer Gradientenabstieg ist rein lokal -- er nutzt keine Information
aus vorherigen Schritten. Die Impulsmethode mischt alte Geschwindigkeit
und aktuellen Gradienten:
\[
  \bm{v}_{\mathrm{neu}} = \beta\,\bm{v}_{\mathrm{alt}}
    - \alpha\,\nabla_{\mathbf{W}} L,
  \qquad 0 \le \beta < 1
\]
$\beta$: Impulsfaktor (inverse Reibung).
 
\medskip
\textbf{Physikalische Interpretation:} Wie eine Kugel, die eine Oberfläche
hinunterrollt und Schwung aufnimmt.
 
\begin{itemize}
  \item I.\,A.\ schneller als klassischer Abstieg.
  \item Lokale Minima können überwunden werden.
  \item \emph{Overshoot:} Kann über das globale Minimum hinausschießen
        $\Rightarrow$ langsames Einpendeln.
\end{itemize}
 
\textbf{Nesterov-Momentum} (beliebt für Mini-Batches):
\[
  \bm{v}_{\mathrm{neu}} = \beta\,\bm{v}_{\mathrm{alt}}
    - \alpha\,\nabla_{\mathbf{W}} L\!\left(\bm{w} + \beta\,\bm{v}_{\mathrm{alt}}\right)
\]
\end{definitionbox}
 
\subsection{Adaptive Lernraten}
 
\begin{definitionbox}{AdaGrad}
Aggregiere die Quadrate der partiellen Ableitungen:
\[
  A_i \leftarrow A_i + \left(\frac{\partial L}{\partial w_i}\right)^2,
  \qquad
  w_i \leftarrow w_i - \frac{\alpha_0}{\sqrt{A_i}}\,\frac{\partial L}{\partial w_i}
\]
\emph{Problem:} historische Fehler werden mitgeschleppt;
zu schnelle Reduktion der Lernrate (``Jugendsünden'').
\end{definitionbox}
 
\begin{definitionbox}{RMSprop}
Exponentielle Mittelung der Gradientquadrate mit $0 < \rho < 1$:
\[
  A_i \leftarrow \rho A_i + (1-\rho)\left(\frac{\partial L}{\partial w_i}\right)^2,
  \qquad
  w_i \leftarrow w_i - \frac{\alpha(t)}{\sqrt{A_i}}\,\frac{\partial L}{\partial w_i}
\]
Historische Probleme werden ``vergessen''. Zusätzliche Lernraten-Abklingsstrategie
$\alpha(t)$ erforderlich.
\end{definitionbox}
 
\begin{importantbox}{Adam -- Adaptive Moment Estimation}
Kombination aus Impulsmethode und adaptiven Lernraten; beim Training
tiefer neuronaler Netze aktuell der verbreitetste Optimierungsalgorithmus:
\[
  A_i \leftarrow \rho A_i + (1-\rho)\left(\frac{\partial L}{\partial w_i}\right)^2
  \quad\text{(wie RMSprop)}
\]
\[
  F_i \leftarrow \rho_f F_i + (1-\rho_f)\,\frac{\partial L}{\partial w_i}
  \quad\text{(Gradientmittelung wie Impuls)}
\]
\[
  w_i \leftarrow w_i - \frac{\tilde{\alpha}_t}{\sqrt{A_i}}\,F_i,
  \quad
  \tilde{\alpha}_t = \alpha_t \cdot\frac{\sqrt{1-\rho^t}}{1-\rho_f^t}
  \quad\text{(Bias-Korrektur)}
\]
In der Praxis wird $\sqrt{A_i}$ durch $\sqrt{A_i + \varepsilon}$
ersetzt (Regularisierung).
\end{importantbox}


% ═══════════════════════════════════════════════════════════════════════════
\section{Frage 10 · Fourier-Reihen und Fourier-Transformation}
% ═══════════════════════════════════════════════════════════════════════════
 
\subsection{Motivation: Signale in Zeit und Frequenz}
 
Viele Signale lassen sich am besten nicht im \emph{Zeitbereich}
(Amplitude als Funktion der Zeit), sondern im \emph{Frequenzbereich}
(Welche Schwingungen stecken im Signal?) beschreiben.
Die Fourier-Theorie liefert den mathematischen Rahmen für diesen Wechsel
der Darstellung.
 
\textbf{Grundidee:} Beliebige (hinreichend gutartige) Funktionen lassen
sich als Überlagerung von Sinus- und Kosinus-Schwingungen (bzw.\
komplexen Exponentialfunktionen) darstellen.
 
\subsection{Orthogonalität -- das Fundament}
 
\begin{definitionbox}{Skalarprodukt für Funktionen}
Für Funktionen $f, g: [-\tfrac{T}{2}, \tfrac{T}{2}] \to \mathbb{R}$
(oder $\mathbb{C}$) definiert man das Skalarprodukt als
\[
  \langle f, g \rangle
  = \intT f(t)\,\overline{g(t)}\,\mathrm{d}t.
\]
Zwei Funktionen heißen \textbf{orthogonal}, wenn $\langle f,g\rangle = 0$.
\end{definitionbox}
 
\begin{importantbox}{Orthogonalität der Fourier-Basisfunktionen}
Die trigonometrischen Funktionen $\{1,\, \cos(n\omega_0 t),\, \sin(n\omega_0 t)\}$
mit Grundkreisfrequenz $\omega_0 = \frac{2\pi}{T}$ bilden ein
\textbf{orthogonales System} auf $[-\tfrac{T}{2}, \tfrac{T}{2}]$:
\[
  \intT \cos(n\omega_0 t)\,\cos(m\omega_0 t)\,\mathrm{d}t
  = \begin{cases} T/2 & n = m \neq 0 \\ 0 & n \neq m \end{cases}
\]
und analog für Sinus; gemischte Produkte $\sin \cdot \cos$ sind stets~$0$.
 
Die komplexen Exponentialfunktionen $\e^{\ii n \omega_0 t}$ sind
ebenfalls orthogonal:
\[
  \intT \e^{\ii n \omega_0 t}\,\overline{\e^{\ii m \omega_0 t}}\,\mathrm{d}t
  = \intT \e^{\ii(n-m)\omega_0 t}\,\mathrm{d}t
  = \begin{cases} T & n = m \\ 0 & n \neq m \end{cases}
\]
\end{importantbox}
 
\subsection{Fourier-Reihen (reell)}
 
\begin{definitionbox}{Fourier-Reihe einer $T$-periodischen Funktion}
Sei $f$ eine $T$-periodische, stückweise stetige Funktion.
Dann gilt (an Stetigkeitsstellen):
\[
  f(t) = \frac{a_0}{2}
  + \sum_{n=1}^{\infty} \Bigl[
      a_n \cos(n\omega_0 t) + b_n \sin(n\omega_0 t)
    \Bigr],
  \qquad \omega_0 = \frac{2\pi}{T}
\]
Die \textbf{Fourier-Koeffizienten} werden durch Ausnutzung der
Orthogonalität berechnet:
\[
  a_n = \frac{2}{T}\intT f(t)\,\cos(n\omega_0 t)\,\mathrm{d}t,
  \qquad
  b_n = \frac{2}{T}\intT f(t)\,\sin(n\omega_0 t)\,\mathrm{d}t
\]
\[
  a_0 = \frac{2}{T}\intT f(t)\,\mathrm{d}t \quad \text{(doppelter Mittelwert)}
\]
\end{definitionbox}
 
\textbf{Anschauung:} Die Koeffizienten $a_n$, $b_n$ geben an, wie viel
von der Schwingung mit Frequenz $\nu_n = n/T$ (bzw.\ $n$-ter Harmonischen)
im Signal steckt. Durch das Skalarprodukt ``projiziert'' man $f$ auf
jede Basisfunktion -- genau wie Koordinaten in der linearen Algebra.
 
\begin{examplebox}{Beispiel: Rechteckschwingung}
Für eine Rechtecksschwingung mit Periode $T$ und Amplitude $1$ erhält man:
\[
  f(t) = \frac{1}{2}
  + \frac{2}{\pi}\sin(\omega_0 t)
  + \frac{2}{3\pi}\sin(3\omega_0 t)
  + \frac{2}{5\pi}\sin(5\omega_0 t)
  + \cdots
  = \frac{1}{2} + \sum_{k=0}^{\infty} \frac{2}{(2k+1)\pi}\sin\!\bigl((2k+1)\omega_0 t\bigr)
\]
Nur ungerade Harmonische tragen bei. Je mehr Terme man summiert,
desto besser wird die Annäherung -- außer an den Sprungstellen.
 
\medskip
\textbf{Gibbs'sches Phänomen:} An Sprungstellen ``überschießt'' die
Fourier-Reihe stets um ca.\ $9\,\%$ der Sprungamplitude, egal wie
viele Terme man nimmt. Die Reihe konvergiert dort gegen das
\emph{arithmetische Mittel} des links- und rechtsseitigen Grenzwerts.
\end{examplebox}
 
\subsection{Komplexe Fourier-Reihen}
 
\begin{definitionbox}{Komplexe Fourier-Reihe}
Mithilfe der Euler-Formel $\e^{\ii\phi} = \cos\phi + \ii\sin\phi$
lässt sich die Fourier-Reihe kompakter schreiben:
\[
  f(t) = \sum_{n=-\infty}^{+\infty} c_n\,\e^{\ii n \omega_0 t}
\]
mit komplexen Koeffizienten
\[
  c_n = \frac{1}{T}\intT f(t)\,\e^{-\ii n \omega_0 t}\,\mathrm{d}t.
\]
Zusammenhang mit reellen Koeffizienten:
\[
  c_0 = \frac{a_0}{2}, \qquad
  c_n = \frac{a_n - \ii b_n}{2}, \qquad
  c_{-n} = \overline{c_n} \quad \text{(für reelles } f\text{)}.
\]
\end{definitionbox}
 
\textbf{Interpretation:} Das Spektrum $\{c_n\}$ enthält dieselbe
Information wie das Zeitsignal, aber in einer anderen Darstellung,
die Frequenzanalyse und Filterung erleichtert.
 
\subsection{Fourier-Transformation (kontinuierlich)}
 
Für \emph{nicht}-periodische Signale auf $\mathbb{R}$ lässt man
$T \to \infty$: aus der diskreten Frequenzfolge $\nu_n = n/T$ wird
ein kontinuierliches Frequenzspektrum.
 
\begin{definitionbox}{Fourier-Transformation}
\[
  \hat{f}(\nu) = \intR f(t)\,\e^{-2\pi\ii \nu t}\,\mathrm{d}t
  \qquad \text{(Hintransformation, Analyse)}
\]
\[
  f(t) = \intR \hat{f}(\nu)\,\e^{+2\pi\ii \nu t}\,\mathrm{d}\nu
  \qquad \text{(Rücktransformation, Synthese)}
\]
$\hat{f}(\nu)$ heißt das \textbf{Spektrum} von $f$. Betrag $|\hat{f}(\nu)|$
gibt die Amplitude, Phase $\arg \hat{f}(\nu)$ die Phasenlage der
Frequenzkomponente $\nu$.
\end{definitionbox}
 
\textbf{Wichtige Eigenschaften:}
\begin{itemize}
  \item \textbf{Linearität:} $\widehat{af+bg} = a\hat{f} + b\hat{g}$.
  \item \textbf{Parsevals Theorem:} Die Gesamtenergie ist im Zeit- und
        Frequenzbereich gleich:
        \[
          \intR |f(t)|^2\,\mathrm{d}t = \intR |\hat{f}(\nu)|^2\,\mathrm{d}\nu.
        \]
  \item \textbf{Faltungssatz:} Faltung im Zeitbereich entspricht
        Multiplikation im Frequenzbereich (und umgekehrt):
        \[
          \widehat{f * g}(\nu) = \hat{f}(\nu)\cdot\hat{g}(\nu).
        \]
  \item \textbf{Verschiebung:} Zeitverschiebung um $t_0$ bewirkt eine
        Phasendrehung im Spektrum:
        $\widehat{f(t-t_0)}(\nu) = \e^{-2\pi\ii\nu t_0}\hat{f}(\nu)$.
\end{itemize}
 
\subsection{Diskrete Fourier-Transformation (DFT)}
 
In der Praxis liegen Signale immer diskret vor. Die DFT überträgt
die Fourier-Theorie auf endliche Folgen.
 
\begin{definitionbox}{DFT und IDFT}
Gegeben eine Folge $(x_0, x_1, \ldots, x_{N-1})$, definiert man:
\[
  X_k = \sum_{n=0}^{N-1} x_n\,\e^{-2\pi\ii kn/N},
  \qquad k = 0, 1, \ldots, N-1
  \qquad \text{(DFT)}
\]
\[
  x_n = \frac{1}{N}\sum_{k=0}^{N-1} X_k\,\e^{+2\pi\ii kn/N},
  \qquad n = 0, 1, \ldots, N-1
  \qquad \text{(IDFT)}
\]
\end{definitionbox}
 
\begin{itemize}
  \item Die $N$ komplexen Koeffizienten $X_k$ enthalten dieselbe
        Information wie die $N$ reellen Abtastwerte $x_n$
        (bei reellen Signalen sind nur $N/2$ Koeffizienten unabhängig
        -- die andere Hälfte ist konjugiert-symmetrisch).
  \item Naiver Aufwand: $\mathcal{O}(N^2)$.
  \item \textbf{Fast Fourier Transform (FFT):} Algorithmus
        (Cooley--Tukey) mit Aufwand $\mathcal{O}(N\log N)$ --
        für $N = 2^p$ besonders effizient.
        In Python: \texttt{numpy.fft.fft}.
\end{itemize}
 
\begin{examplebox}{Anwendung: JPG-Bildkompression}
JPEG verwendet eine \textbf{2D-Diskrete-Kosinus-Transformation}
(eng verwandt mit der DFT) auf $8\times 8$-Pixel-Blöcken.
Nur gut sichtbare (niederfrequente) Koeffizienten werden gespeichert --
Hochfrequenzanteile werden verworfen. An harten Kanten entstehen
dabei sichtbare \textbf{Gibbs'sche Spitzen} (``JPEG-Artefakte'').
\end{examplebox}
 
% ═══════════════════════════════════════════════════════════════════════════
\section{Frage 11 · Aliasing-Effekt und Shannon-Nyquist-Theorem}
% ═══════════════════════════════════════════════════════════════════════════
 
\subsection{Kontinuierliche vs.\ diskrete Signale}
 
In der Praxis müssen kontinuierliche Signale \textbf{diskretisiert}
werden, um sie digital verarbeiten zu können. Dabei gibt es zwei Arten
der Diskretisierung:
 
\begin{itemize}
  \item \textbf{Wertdiskretisierung} (Quantisierung): Das Signal wird
        auf endlich viele Amplitudenstufen gerundet.
        Meist weniger kritisch (z.\,B.\ 8-bit vs.\ 16-bit Musik --
        Unterschied in der Dynamik).
  \item \textbf{Zeitliche (räumliche) Diskretisierung} (Abtastung):
        Das Signal wird nur zu bestimmten Zeitpunkten $t_k = k\,\Delta t$
        gemessen. Dies kann zu \emph{extremen Fehlschlüssen} führen.
\end{itemize}
 
\subsection{Der Aliasing-Effekt}
 
\begin{definitionbox}{Aliasing}
Wird ein Signal mit einer Frequenz $\nu$ abgetastet, die höher ist als
die halbe Abtastfrequenz, so kann die Schwingung nicht korrekt
rekonstruiert werden. Sie wird stattdessen als eine andere (niedrigere)
Frequenz $\nu_{\mathrm{alias}}$ wahrgenommen -- eine \textbf{Alias-Frequenz}.
 
\medskip
\textbf{Anschauung:} Zwei verschiedene Sinusschwingungen mit
unterschiedlichen Frequenzen $\nu$ und $\nu_{\mathrm{alias}}$
erzeugen an den Abtastzeitpunkten $t_k = k\,\Delta t$
exakt dieselben Messwerte -- sie sind damit für das diskrete System
\emph{ununterscheidbar}.
\end{definitionbox}
 
\begin{examplebox}{Stroboskop-Effekt}
Ein Wagenrad, das sich schnell dreht, wirkt im Film manchmal
rückwärts drehend oder steht still -- weil die Bildwiederholrate
(Abtastfrequenz) geringer ist als die doppelte Umdrehungsfrequenz.
Ähnlich sieht man bei zu grob abgetasteten Sinussignalen eine
langsamere Schwingung, die es in Wirklichkeit nicht gibt.
\end{examplebox}
 
\subsection{Shannon-Nyquist-Theorem}
 
\begin{importantbox}{Shannon-Nyquist-Abtasttheorem}
Bei äquidistanter Abtastung eines Signals mit Zeitschritt $\Delta t$
(entsprechend einer \textbf{Abtastfrequenz} $\nu_s = 1/\Delta t$)
können nur Frequenzen
\[
  \nu < \nu_{\mathrm{Nyquist}} = \frac{1}{2\,\Delta t} = \frac{\nu_s}{2}
\]
korrekt erkannt und rekonstruiert werden.
 
Höhere Frequenzen werden \textbf{maskiert}, d.\,h.\ auf niedrigere
Alias-Frequenzen abgebildet:
\[
  \nu_{\mathrm{alias}} = \left|\nu - k\,\nu_s\right|,
  \quad k \in \mathbb{Z} \text{ so gewählt, dass }
  \nu_{\mathrm{alias}} \in [0, \nu_{\mathrm{Nyquist}}].
\]
\end{importantbox}
 
\textbf{Beispiele:}
\begin{itemize}
  \item Audio-CD: Abtastrate $44{,}1\,\mathrm{kHz}$
        $\Rightarrow$ Nyquist-Frequenz $22{,}05\,\mathrm{kHz}$
        (deckt das menschliche Hörvermögen bis $\approx 20\,\mathrm{kHz}$ ab).
  \item Medizin/EKG: Das Nyquist-Kriterium muss für alle relevanten
        Herzfrequenzkomponenten erfüllt sein.
  \item Meteorologie: Zu grob aufgelöste Wettermodelle übersehen
        kleinräumige, aber relevante Phänomene.
\end{itemize}
 
\subsection{Konsequenzen und Gegenmaßnahmen}
 
\begin{warnbox}{Achtung: Aliasing ist nicht reparierbar!}
Ist das Signal bereits abgetastet und Aliasing aufgetreten, lassen sich
die ursprünglichen hohen Frequenzen \emph{nicht mehr} rekonstruieren --
die Information ist unwiederbringlich verloren. Aliasing muss daher
\emph{vor} der Abtastung verhindert werden.
\end{warnbox}
 
\textbf{Gegenmaßnahmen:}
\begin{enumerate}
  \item \textbf{Höhere Abtastrate:} $\Delta t$ verkleinern, sodass
        $\nu_{\mathrm{Nyquist}}$ alle im Signal vorkommenden Frequenzen
        übersteigt. Oft teuer (mehr Speicher, schnellere Hardware).
  \item \textbf{Anti-Aliasing-Filter:} Vor der Abtastung ein
        \textbf{Tiefpassfilter} anwenden, der alle Frequenzen oberhalb
        von $\nu_{\mathrm{Nyquist}}$ herausfiltert. So ist sichergestellt,
        dass keine Frequenzen abgetastet werden, die das Theorem verletzen.
\end{enumerate}
 
\subsection{Bedeutung von Domänenwissen}
 
Das Shannon-Nyquist-Theorem macht deutlich, warum \textbf{Domänenwissen}
bei der Datenerhebung unerlässlich ist: Man muss \emph{vorab} wissen,
welche Frequenzen im zu messenden Prozess physikalisch relevant sind,
um die Abtastrate sinnvoll zu wählen.
 
\begin{itemize}
  \item Zu grobe Abtastung $\Rightarrow$ Aliasing, falsche Schlussfolgerungen.
  \item Zu feine Abtastung $\Rightarrow$ unnötig große Datenmengen,
        höhere Kosten, langsamere Verarbeitung.
  \item Optimale Wahl erfordert Kenntnis des Prozesses.
\end{itemize}
 
\begin{examplebox}{Spektralanalyse mit der DFT}
Gegeben ein Signal mit $N = 256$ Abtastwerten und Zeitschritt $\Delta t$.
 
Die DFT liefert $N$ Koeffizienten $X_k$. Die dabei auftretenden
\textbf{Frequenzachse} reicht von $0$ bis $\nu_s = 1/\Delta t$,
aufgeteilt in $N$ Schritte. Wegen der konjugiert-symmetrischen
Natur der DFT (bei reellen Signalen) enthält nur die erste Hälfte
($k = 0, \ldots, N/2$) unabhängige Information -- bis zur
Nyquist-Frequenz $\nu_{\mathrm{Nyquist}} = 1/(2\Delta t)$.
 
In den $N/2 = 128$ komplexen Koeffizienten steckt dieselbe Information
wie in den $N = 256$ reellen Abtastwerten -- allerdings im Frequenzbereich
dargestellt, was andere Bearbeitungsmöglichkeiten eröffnet
(Filterung, Kompression, Mustererkennung).
\end{examplebox}
 
\subsection{Zusammenfassung: Zeit- vs.\ Frequenzbereich}
 
\begin{center}
\begin{tabular}{lll}
  \toprule
  & \textbf{Zeitbereich} & \textbf{Frequenzbereich} \\
  \midrule
  Darstellung & $f(t)$ bzw.\ $x_n$ & $\hat{f}(\nu)$ bzw.\ $X_k$ \\
  Interpretation & Signal als Zeitverlauf & Spektrum der Schwingungen \\
  Faltung & $f * g$ & $\hat{f} \cdot \hat{g}$ (Multiplikation) \\
  Filterung & aufwändig & einfach (Koeffizienten nullsetzen) \\
  Übergang & \multicolumn{2}{c}{$\longleftrightarrow$ Fourier-Transformation / DFT / FFT} \\
  \bottomrule
\end{tabular}
\end{center}


% ═══════════════════════════════════════════════════════════════════════════
\section{Frage 12 · Faltung (diskret und stetig) sowie Faltungskerne und Anwendungen}
% ═══════════════════════════════════════════════════════════════════════════
 
\subsection{Motivation und Grundidee}
 
Die \textbf{Faltung} beschreibt, wie ein Signal durch ein System
verändert wird -- oder allgemeiner, wie zwei Funktionen miteinander
``verschmiert'' werden. Sie tritt überall dort auf, wo ein System
linear und zeitinvariant (LTI) ist: Audio-Effekte, Bildverarbeitung,
neuronale Netze, Differenzialgleichungen.
 
\textbf{Kernidee:} Der Ausgabewert zum Zeitpunkt $t$ hängt nicht nur
vom aktuellen Eingabewert ab, sondern von einem \emph{gewichteten
Mittel} der Vergangenheit (oder Umgebung) -- das Gewicht gibt der
Faltungskern (auch: Filter, Impulsantwort) vor.
 
\subsection{Stetige Faltung}
 
\begin{definitionbox}{Stetige Faltung}
Für zwei Funktionen $f, g: \mathbb{R} \to \mathbb{R}$ ist die
Faltung definiert als
\[
  (f * g)(t)
  = \intR f(\tau)\,g(t - \tau)\,\mathrm{d}\tau
  = \intR f(t - \tau)\,g(\tau)\,\mathrm{d}\tau.
\]
$g$ heißt \textbf{Faltungskern} (oder Impulsantwort, Filter).
 
\medskip
\textbf{Eigenschaften:}
\begin{itemize}
  \item \emph{Kommutativität:} $f * g = g * f$
  \item \emph{Assoziativität:} $(f * g) * h = f * (g * h)$
  \item \emph{Linearität:} $(af + bg) * h = a(f*h) + b(g*h)$
  \item \emph{Faltungssatz:} Faltung im Zeitbereich $\;\widehat{=}\;$
        Multiplikation im Frequenzbereich:
        \[
          \widehat{f * g}(\nu) = \hat{f}(\nu)\cdot\hat{g}(\nu)
        \]
\end{itemize}
\end{definitionbox}
 
\subsection{Diskrete Faltung}
 
In der Praxis arbeitet man mit diskreten Signalen $(x_n)$ und einem
diskreten Kern $(k_m)$.
 
\begin{definitionbox}{Diskrete Faltung}
\[
  (x * k)_n = \sum_{m=-\infty}^{+\infty} x_m\,k_{n-m}
  = \sum_{m=-\infty}^{+\infty} x_{n-m}\,k_m
\]
Bei einem endlichen Kern der Länge $M$ (z.\,B.\ $m = 0, \ldots, M-1$)
vereinfacht sich die Summe auf $M$ Terme pro Ausgabewert.
 
\medskip
\textbf{Randbehandlung} (bei endlichen Signalen):
\begin{itemize}
  \item \emph{Zero-padding:} Signal außerhalb mit $0$ auffüllen.
  \item \emph{Periodische Fortsetzung:} Signal wird zyklisch wiederholt.
  \item \emph{Spiegelung:} Randwerte werden gespiegelt.
\end{itemize}
Die Wahl beeinflusst die Ergebnisse an den Rändern.
\end{definitionbox}
 
\begin{examplebox}{Rechenbeispiel: Diskrete Faltung}
Signal $x = (1, 2, 3, 4)$, Kern $k = (1, 0, -1)$ (Differenzoperator):
\[
  (x * k)_n = x_n \cdot 1 + x_{n-1} \cdot 0 + x_{n-2} \cdot (-1)
  = x_n - x_{n-2}
\]
Ergebnis (mit Zero-padding): $(1,\; 2,\; 2,\; 2,\; -3,\; -4)$.
Dieser Kern approximiert eine erste Ableitung -- er ist ein einfacher
\textbf{Differenzierungskern}.
\end{examplebox}
 
\subsection{Wichtige Faltungskerne und ihre Anwendungen}
 
\begin{definitionbox}{Übersicht: Faltungskerne}
\renewcommand{\arraystretch}{1.4}
\begin{tabular}{p{3.5cm} p{4cm} p{5.5cm}}
  \toprule
  \textbf{Kern} & \textbf{Formel (diskret)} & \textbf{Wirkung / Anwendung} \\
  \midrule
  Glättung / Mittelwert &
  $k = \tfrac{1}{M}(1, 1, \ldots, 1)$ &
  Rauschunterdrückung; Tiefpassfilter \\
  Gauß-Kern &
  $k_m \propto \mathrm{e}^{-m^2/(2\sigma^2)}$ &
  Weiche Glättung; kein Gibbs-Effekt \\
  Differenzierungskern &
  $k = (1, -1)$ oder $(-1, 0, 1)$ &
  Kantendetektion (Hochpassfilter) \\
  Laplace-Kern (2D) &
  $\begin{pmatrix}0&1&0\\1&-4&1\\0&1&0\end{pmatrix}$ &
  Kantendetektion in Bildern \\
  Sobel-Kern (2D) &
  $\begin{pmatrix}-1&0&1\\-2&0&2\\-1&0&1\end{pmatrix}$ &
  Gerichtete Kantendetektion ($x$-Richtung) \\
  Dirac-Delta &
  $k = (\ldots, 0, 1, 0, \ldots)$ &
  Identität: $f * \delta = f$ \\
  \bottomrule
\end{tabular}
\end{definitionbox}
 
\textbf{Tiefpass vs.\ Hochpass:}
\begin{itemize}
  \item \textbf{Tiefpassfilter} (z.\,B.\ Gauß, Mittelwert): dämpfen
        hohe Frequenzen, glätten das Signal $\Rightarrow$ Rauschreduktion.
  \item \textbf{Hochpassfilter} (z.\,B.\ Differenz, Laplace): dämpfen
        niedrige Frequenzen, betonen Änderungen $\Rightarrow$ Kantendetektion.
\end{itemize}
 
\subsection{Anwendungen der Faltung}
 
\textbf{Signalverarbeitung:}
\begin{itemize}
  \item Tiefpass-Filterung von Messsignalen (Rauschunterdrückung).
  \item Hall-Effekt in der Audiotechnik: Raumimpulsantwort wird mit dem
        Trockensignal gefaltet.
  \item Gleitender Mittelwert (Moving Average) in Zeitreihen.
\end{itemize}
 
\textbf{Bildverarbeitung:}
\begin{itemize}
  \item Weichzeichnen (Gauß-Blur), Schärfen, Kantendetektion (Sobel, Laplace).
  \item JPEG-Kompression nutzt die enge Verwandtschaft zwischen
        Faltung und Kosinus-Transformation.
\end{itemize}
 
\textbf{Convolutional Neural Networks (CNNs):}
\begin{itemize}
  \item Die Kerngewichte sind \emph{lernbare} Faltungskerne.
  \item Durch Training lernt das Netz, welche lokalen Muster
        (Kanten, Texturen, Formen) für die Aufgabe relevant sind.
  \item Translationsinvarianz: derselbe Kern wird über das gesamte
        Bild geschoben $\Rightarrow$ enorme Parameterreduktion
        gegenüber vollständig vernetzten Schichten.
\end{itemize}
 
\begin{importantbox}{Faltung und Fourier-Transformation}
Der \textbf{Faltungssatz} macht die enge Verbindung zur Fourier-Theorie
deutlich: Eine Faltung im Zeitbereich (rechenintensiv: $\mathcal{O}(N^2)$)
entspricht einer Multiplikation im Frequenzbereich. Mithilfe der FFT
lässt sich die Faltung daher effizient in $\mathcal{O}(N \log N)$
berechnen:
\[
  x * k = \mathcal{F}^{-1}\!\bigl(\mathcal{F}(x) \cdot \mathcal{F}(k)\bigr)
\]
\end{importantbox}
 
% ═══════════════════════════════════════════════════════════════════════════
\section{Frage 13 · Methoden zur numerischen Bestimmung von Integralen}
% ═══════════════════════════════════════════════════════════════════════════
 
\subsection{Motivation}
 
\begin{itemize}
  \item Viele naturwissenschaftlich-technische Aufgaben sind mittels
        Integration lösbar (Flächen, Volumen, Energien, Wahrscheinlichkeiten).
  \item Auswertung statistischer Größen führt sehr oft auf Integration
        (z.\,B.\ Erwartungswert $\mathbb{E}[X] = \int x\,p(x)\,\mathrm{d}x$).
  \item \textbf{Symbolisches Integrieren} ist meist schwierig oder gar
        nicht möglich.
  \item \textbf{Numerisches Integrieren} ist im Vergleich zum numerischen
        Differenzieren eine \emph{relativ gutartige} Aufgabe --
        zumindest für ein- und niedrigdimensionale Integrale.
\end{itemize}
 
\subsection{Grundprinzip: Newton-Cotes-Formeln}
 
\begin{definitionbox}{Allgemeines Schema (normiert)}
Das Integral wird durch eine gewichtete Summe von Funktionswerten
an äquidistanten Stützstellen approximiert:
\[
  \int_a^b f(x)\,\mathrm{d}x
  \;\approx\; (b-a)\sum_{k=0}^{n} w_k\, f(x_k)
\]
wobei $x_k = a + k\,h$ mit $h = \frac{b-a}{n}$ und die Gewichte
$w_k$ die jeweilige Formel charakterisieren.
 
\medskip
\textbf{Forderungen an die Gewichte:}
\begin{itemize}
  \item $\sum_{k=0}^{n} w_k = 1$ (sonst wäre das Ergebnis für $f \equiv 1$ falsch).
  \item Meistens auch $w_k > 0$ (sonst droht Auslöschung bei negativen Gewichten).
\end{itemize}
\end{definitionbox}
 
\subsection{Trapezregel}
 
\begin{definitionbox}{Trapezregel (summiert)}
Jedes Teilintervall $[x_k, x_{k+1}]$ wird durch ein Trapez approximiert
(lineare Interpolation zwischen den Endpunkten):
\[
  \int_a^b f(x)\,\mathrm{d}x
  \approx \frac{h}{2}\Bigl(f(x_0) + 2\sum_{k=1}^{n-1} f(x_k) + f(x_{n})\Bigr),
  \qquad h = \frac{b-a}{n}
\]
Gewichte (normiert): $w_0 = w_n = \tfrac{1}{2(n)}$,
$w_k = \tfrac{1}{n}$ für $1 \le k \le n-1$.
 
\medskip
\textbf{Fehlerordnung:} $\mathcal{O}(h^2)$ pro Teilintervall,
$\mathcal{O}(h^2)$ global (die Einzelfehler summieren sich, aber
ihre Anzahl wächst auch mit $1/h$).
\end{definitionbox}
 
\subsection{Simpson-Regel}
 
\begin{definitionbox}{Simpson-Regel (summiert)}
Statt linearer Interpolation wird auf je \emph{drei} Stützpunkten
ein Parabel-Stück gelegt (quadratische Interpolation).
Für gerades $n$:
\[
  \int_a^b f(x)\,\mathrm{d}x
  \approx \frac{h}{3}\Bigl(
    f(x_0)
    + 4\sum_{\substack{k=1\\k\,\text{ungerade}}}^{n-1} f(x_k)
    + 2\sum_{\substack{k=2\\k\,\text{gerade}}}^{n-2} f(x_k)
    + f(x_n)
  \Bigr)
\]
Gewichte für 2 Teilintervalle (Bsp.): $\tfrac{1}{6},\, \tfrac{4}{6},\, \tfrac{1}{6}$.
Für allgemeine Newton-Cotes-Formeln mit mehr Punkten:
$\tfrac{1}{12},\, \tfrac{3}{12},\, \tfrac{6}{12},\, \tfrac{3}{12},\, \tfrac{1}{12}$ (4 Intervalle).
 
\medskip
\textbf{Fehlerordnung:} $\mathcal{O}(h^4)$ -- deutlich besser als die Trapezregel,
da die Parabel bereits kubische Polynome exakt integriert
(\emph{Genauigkeitsgrad} 3).
\end{definitionbox}
 
\begin{examplebox}{Vergleich Trapez vs.\ Simpson}
Für $\int_0^1 \sin(\pi x)\,\mathrm{d}x = \tfrac{2}{\pi} \approx 0{,}6366$
mit $n = 4$ Teilintervallen ($h = 0{,}25$):
\begin{itemize}
  \item Trapezregel: $\approx 0{,}6088$ \quad (Fehler $\approx 4{,}4\%$)
  \item Simpson-Regel: $\approx 0{,}6366$ \quad (Fehler $< 0{,}01\%$)
\end{itemize}
\end{examplebox}
 
\subsection{Romberg-Integration}
 
\begin{definitionbox}{Romberg-Integration}
Idee: Kombiniere Trapezregeln mit unterschiedlicher Schrittweite
($n = 1, 2, 4, 8, \ldots$ Teilintervalle) zu einer Folge verbesserter
Näherungen durch \textbf{Richardson-Extrapolation}.
 
Bezeichnet man die Trapezregel mit $n$ Stützstellen und Ordnung $m$
als $Q_n^{(m)}$, so ergibt sich die verbesserte Näherung:
\[
  Q_n^{(m+1)} = Q_n^{(m)} + \frac{Q_n^{(m)} - Q_{n/2}^{(m)}}{4^m - 1}
\]
Das \textbf{Romberg-Schema} berechnet so ein Dreieck von Näherungen:
\[
  Q_1^{(1)} \to Q_1^{(2)} \to Q_1^{(3)} \to \cdots
\]
\[
  Q_2^{(1)} \nearrow Q_2^{(2)} \nearrow \cdots
\]
\[
  Q_4^{(1)} \nearrow \cdots
\]
Jede Spalte erhöht die Fehlerordnung um $\mathcal{O}(h^2)$;
das Verfahren konvergiert oft sehr rasch.
 
\medskip
\textbf{Vorteil:} Weiterhin äquidistante Stützstellen;
bereits berechnete Funktionswerte werden wiederverwendet.
\end{definitionbox}
 
\subsection{Gauß-Legendre-Integration}
 
\begin{definitionbox}{Gauß-Legendre-Integration}
\textbf{Grundidee:} Stützstellen sind \emph{nicht} mehr äquidistant,
sondern werden \emph{als zusätzliche Freiheitsgrade} optimal gewählt.
 
Mit $n$ Stützstellen $x_1, \ldots, x_n$ und Gewichten
$w_1, \ldots, w_n$ auf $[-1, 1]$:
\[
  \int_{-1}^{1} f(x)\,\mathrm{d}x \approx \sum_{k=1}^{n} w_k\,f(x_k)
\]
Die optimale Wahl (Nullstellen der Legendre-Polynome) ergibt einen
\textbf{Genauigkeitsgrad von $2n-1$}: Polynome bis Grad $2n-1$
werden exakt integriert -- mit nur $n$ Funktionsauswertungen!
 
\medskip
\textbf{Beispiel ($n = 2$):}
\[
  x_{1,2} = \pm\tfrac{1}{\sqrt{3}}, \qquad w_{1,2} = 1
  \qquad\Rightarrow\quad \text{exakt für Polynome bis Grad 3}
\]
 
\textbf{Vorteil} gegenüber Newton-Cotes: Bei gleicher Anzahl von
Funktionsauswertungen deutlich höhere Genauigkeit.
 
\textbf{Nachteil:} Die Stützstellen sind i.\,A.\ keine ``schönen''
Zahlen und fest vorgegeben; das Verfahren eignet sich weniger gut,
wenn Funktionswerte nur an äquidistanten Stellen vorliegen.
\end{definitionbox}
 
\subsection{Übersicht und Vergleich der Methoden}
 
\begin{center}
\renewcommand{\arraystretch}{1.35}
\begin{tabular}{lccll}
  \toprule
  \textbf{Methode} & \textbf{Stützst.} & \textbf{Äquidist.?}
    & \textbf{Fehlerordnung} & \textbf{Besonderheit} \\
  \midrule
  Trapezregel     & $n+1$ & ja  & $\mathcal{O}(h^2)$   & einfach, robust \\
  Simpson-Regel   & $n+1$ & ja  & $\mathcal{O}(h^4)$   & $n$ gerade nötig \\
  Romberg         & variabel & ja & sehr hoch (iterativ) & Richardson-Extrapolation \\
  Gauß-Legendre   & $n$   & nein & $\mathcal{O}(h^{2n})$ & opt.\ Stützstellen \\
  Monte Carlo     & $N$   & zufällig & $\mathcal{O}(N^{-1/2})$ & hochdimensional \\
  \bottomrule
\end{tabular}
\end{center}
 
\subsection{Hochdimensionale Integrale: Monte-Carlo-Integration}
 
Für Integrale in hohen Dimensionen versagen alle oben genannten
Methoden -- der Aufwand wächst \emph{exponentiell} mit der Dimension
(\emph{Fluch der Dimensionalität}): Für $d$ Dimensionen und $n$
Stützstellen pro Achse benötigt man $n^d$ Funktionsauswertungen.
 
\begin{definitionbox}{Monte-Carlo-Integration}
\[
  \int_\Omega f(x)\,\mathrm{d}x \approx \frac{|\Omega|}{N}\sum_{i=1}^{N} f(x_i),
  \qquad x_i \sim \mathcal{U}(\Omega)
\]
Die Stützstellen werden \textbf{zufällig} gleichverteilt in $\Omega$
gezogen. Der Fehler skaliert mit $\mathcal{O}(N^{-1/2})$ --
\emph{unabhängig von der Dimension}. Das macht Monte-Carlo zur
Standardmethode für hochdimensionale Integrale.
\end{definitionbox}
 
\begin{warnbox}{Numerisches Differenzieren vs.\ Integrieren}
Numerisches \textbf{Integrieren} ist eine \emph{gutartige} Operation:
Integration ``glättet'' -- kleine Fehler in $f$ werden durch die
Summation herausgemittelt.
 
Numerisches \textbf{Differenzieren} ist \emph{heikel}: Differenzierung
wirkt als \textbf{Hochpassfilter} und verstärkt Rauschen. Je größer
die Frequenz $\omega$, desto stärker die Verstärkung
($\frac{d}{dt}\sin(\omega t) = \omega\cos(\omega t)$).
Daher soll man höhere Ableitungen nicht durch mehrfaches numerisches
Differenzieren berechnen, sondern direkt mit einer analytisch
hergeleiteten Formel (z.\,B.\ via Taylor-Entwicklung):
\[
  f''(x_0) \approx \frac{f(x_0+h) - 2f(x_0) + f(x_0-h)}{h^2}
\]
\end{warnbox}











% ─────────────────────────────────────────────────────────────────────────────
\subsection{Grundbegriffe der Graphentheorie}

\definition{Ein \textbf{Graph} $G = (V, E)$ besteht aus einer \textbf{Knotenmenge} $V$ (engl.\ \emph{vertices}) und einer \textbf{Kantenmenge} $E$ (engl.\ \emph{edges}). Jeder Kante $e \in E$ sind zwei (nicht notwendigerweise verschiedene) Knoten zugeordnet, die als \emph{inzident} zu $e$ bezeichnet werden. Graphen dienen der mathematischen Modellierung netzartiger Strukturen; die konkrete Beschaffenheit der Knoten und Kanten spielt dabei keine Rolle (topologische Abstraktion).}

\textbf{Grundbegriffe:}
\begin{itemize}[leftmargin=1.5em]
    \item \textbf{Schlinge:} Eine Kante der Form $e = \{v, v\}$ (Endknoten identisch). Schlingen erschweren viele graphentheoretische Analysen erheblich.
    \item \textbf{Parallele Kanten:} Zwei Kanten $a = \{v,w\}$ und $b = \{v,w\}$ mit identischen Endknoten. Graphen ohne Schlingen und parallele Kanten heißen \textbf{schlicht} (einfach); andernfalls \textbf{Multigraph}.
    \item \textbf{Nachbarschaft:} $\Gamma(v) = \{w \in V \mid \{v,w\} \in E,\, w \neq v\}$ - die Menge aller zu $v$ benachbarten Knoten. Ein Knoten mit $\Gamma(v) = \emptyset$ heißt \textbf{isoliert}.
    \item \textbf{Grad:} $\deg(v)$ = Anzahl der von $v$ ausgehenden Kanten (Schlingen zählen doppelt). Für schlichte Graphen gilt $\deg(v) = |\Gamma(v)|$. Es gilt stets $\sum_{v \in V} \deg(v) = 2|E|$ (Handshaking-Lemma), d.\,h.\ die Anzahl der Knoten mit ungeradem Grad ist stets gerade.
    \item \textbf{Maximalgrad / Minimalgrad / Mittlerer Grad:} $\Delta(G) = \max_{v}\deg(v)$, $\delta(G) = \min_{v}\deg(v)$, $\Phi(G) = \frac{2|E|}{|V|}$.
    \item \textbf{Dichte} (schlichte Graphen): $\rho(G) = \frac{|E|}{\binom{|V|}{2}} = \frac{2|E|}{|V|(|V|-1)}$ - Verhältnis der vorhandenen Kanten zur maximalen Kantenzahl.
\end{itemize}

\textbf{Wege und Abstände:}
\begin{itemize}[leftmargin=1.5em]
    \item \textbf{Kantenfolge:} Folge $v_1, e_1, v_2, e_2, \ldots, v_{k-1}, e_{k-1}, v_k$, wobei $e_i = \{v_i, v_{i+1}\}$. Die \textbf{Länge} ist die Anzahl der Kanten.
    \item \textbf{Weg:} Kantenfolge, in der jeder Knoten höchstens einmal auftritt. (Kanten dürfen bei einem Weg nicht doppelt auftreten, da sonst ein bereits besuchter Knoten erneut besucht würde.)
    \item \textbf{Kreis:} Geschlossene Kantenfolge, bei der erster und letzter Knoten identisch sind. Die Länge des kürzesten Kreises heißt \textbf{Taille}, die des längsten \textbf{Umfang}. Eine Kante, die auf keinem Kreis liegt, heißt \textbf{Brücke}.
    \item \textbf{Abstand:} $\mathrm{dist}(v,w)$ = Länge des kürzesten Weges zwischen $v$ und $w$; $\mathrm{dist}(v,v) = 0$; $\mathrm{dist}(v,w) = \infty$ falls kein Weg existiert.
    \item \textbf{Exzentrizität:} $e(v) = \max_{w \in V} \mathrm{dist}(v,w)$ - Abstand zum am weitesten entfernten Knoten.
    \item \textbf{Durchmesser, Radius, Zentrum:} $d(G) = \max_v e(v)$; $r(G) = \min_v e(v)$; Knoten mit $e(v) = r(G)$ liegen im \textbf{Zentrum}. Es gilt $r(G) \leq d(G) \leq 2r(G)$.
\end{itemize}

\textbf{Zusammenhang:} Ein Graph heißt \textbf{zusammenhängend}, wenn zwischen allen Knotenpaaren ein Weg existiert (kein Abstand ist $\infty$). Andernfalls zerfällt er in mehrere \textbf{Komponenten}. Äquivalent: $G$ ist zusammenhängend, wenn er durch sukzessive Kantenkontraktionen auf einen einzelnen Knoten reduziert werden kann.

\newpage
% ─────────────────────────────────────────────────────────────────────────────
\subsection{Graphentheorie: Inzidenz-, Grad-, Adjazenz-, Laplace- und Abstandsmatrizen sowie deren Zusammenhang}

Sei $G = (V,E)$ ein Graph mit $n = |V|$ Knoten und $m = |E|$ Kanten.

\textbf{Inzidenzmatrix $B(G)$} ($n \times m$): Eintrag $b_{ve} = 1$, wenn Knoten $v$ und Kante $e$ inzidieren; bei Schlingen wird $b_{ve} = 2$ gesetzt. Jede Spaltensumme ist 2 (jede Kante hat genau zwei Endknoten).

\textbf{Gradmatrix $D(G)$} ($n \times n$): Diagonalmatrix $D = \mathrm{diag}(\deg(1), \deg(2), \ldots, \deg(n))$. Die Diagonaleinträge entsprechen den Knotengraden.

\textbf{Adjazenzmatrix $A(G)$} ($n \times n$): Eintrag $a_{ij}$ = Anzahl der Kanten zwischen Knoten $i$ und $j$ (Schlingen je nach Konvention einfach oder doppelt). Wichtige Eigenschaften:
\begin{itemize}[leftmargin=1.5em]
    \item $A(G)$ ist stets \textbf{symmetrisch}. Bei schlichten Graphen enthält sie nur 0 und 1, mit Nullen auf der Diagonale.
    \item Die Zeilen-/Spaltensumme entspricht dem Grad des jeweiligen Knotens (wenn Schlingen doppelt gezählt werden).
    \item \textbf{Potenzen der Adjazenzmatrix:} Das Element $(A^p)_{ij}$ gibt die Anzahl aller Kantenfolgen der Länge $p$ von Knoten $i$ nach Knoten $j$ an. Damit lassen sich Abstände und Erreichbarkeit systematisch ermitteln.
    \item Lässt sich $A(G)$ durch Permutation in eine Blockstruktur $\begin{pmatrix} X & 0 \\ 0 & Y \end{pmatrix}$ überführen, ist der Graph \textbf{nicht zusammenhängend}.
\end{itemize}

\textbf{Laplace-Matrix $L(G)$} und \textbf{vorzeichenlose Laplace-Matrix $Q(G)$}:
\formel{L(G) = D(G) - A(G), \qquad Q(G) = D(G) + A(G) = B(G)\cdot B(G)^\top}
Die Laplace-Matrix ist positiv semidefinit, hat stets mindestens einen Eigenwert 0, und die Anzahl der Eigenwerte 0 entspricht der Anzahl der Zusammenhangskomponenten.

\textbf{Abstandsmatrix $\mathcal{M}(G)$} ($n \times n$): Eintrag $m_{ij} = \mathrm{dist}(i,j)$. Berechnung durch sukzessives Auswerten der Potenzen der Adjazenzmatrix: Sobald $(A^p)_{ij} > 0$ für ein noch nicht eingetragenes Paar, setzt man $m_{ij} = p$.

\textbf{Zusammenhang der Matrizen:}
\begin{itemize}[leftmargin=1.5em]
    \item $\|B(G)\|_{H1} = \|D(G)\|_{H1} = \|A(G)\|_{H1} = 2|E|$ (Summe aller Einträge, bei doppelter Schlinge)
    \item $Q(G) = B(G) \cdot B(G)^\top = D(G) + A(G)$
    \item $L(G) = D(G) - A(G)$
    \item Exzentrizitäten, Durchmesser und Radius lassen sich direkt aus $\mathcal{M}(G)$ ablesen.
\end{itemize}

\newpage
% ─────────────────────────────────────────────────────────────────────────────
\subsection{Graphentheorie: Eulersche Kantenfolgen und Eulersche Touren sowie Hamiltonwege und Hamiltonkreise}

\textbf{Definitionen:}
\begin{itemize}[leftmargin=1.5em]
    \item \textbf{Eulersche Kantenfolge} (Eulerscher Weg): Eine Kantenfolge, die jede Kante von $G$ \emph{genau einmal} durchläuft. Im Allgemeinen kein Weg, da Knoten mehrfach besucht werden dürfen.
    \item \textbf{Eulersche Tour} (Eulerkreis): Eine \emph{geschlossene} Kantenfolge, die jede Kante genau einmal durchläuft. Graphen mit einer Eulerschen Tour heißen \textbf{eulersch}.
    \item \textbf{Hamiltonweg}: Ein \emph{Weg} in $G$, der jeden Knoten aus $V$ genau einmal enthält.
    \item \textbf{Hamiltonkreis}: Ein \emph{Kreis} in $G$, der jeden Knoten aus $V$ enthält. Graphen mit einem Hamiltonkreis heißen \textbf{hamiltonsch}.
\end{itemize}

\textbf{Existenzsätze für Eulersche Strukturen:}

\definition{\textbf{Satz (Euler-Hierholzer):} Ein zusammenhängender Graph enthält genau dann eine \textbf{Eulersche Tour}, wenn \emph{jeder} Knoten einen \textbf{geraden Grad} hat ($\forall v \in V: \deg(v) \bmod 2 = 0$). Er enthält genau dann eine \textbf{Eulersche Kantenfolge}, wenn \emph{maximal zwei} Knoten einen ungeraden Grad haben.}

\beispiel{Das \textbf{Königsberger Brückenproblem} (Euler, 1736): Sieben Brücken verbinden Stadtteile von Königsberg. Alle vier Knoten des entsprechenden Graphen haben ungeraden Grad $\Rightarrow$ keine Eulersche Kantenfolge existiert.}

\textbf{Existenzsätze für Hamiltonsche Strukturen:}

\definition{\textbf{Satz von Dirac:} Sei $G$ ein schlichter Graph mit $|V| \geq 3$ und $\deg(v) \geq \frac{|V|}{2}$ für alle $v \in V$. Dann besitzt $G$ einen Hamiltonkreis. (\emph{hinreichend, nicht notwendig})}

\definition{\textbf{Satz von Ore:} Sei $G$ schlicht mit $|V| \geq 3$. Gilt für alle nicht-benachbarten Paare $u,v \in V$: $\deg(u) + \deg(v) \geq |V|$, so besitzt $G$ einen Hamiltonkreis. (\emph{hinreichend, nicht notwendig; verallgemeinert Dirac})}

\textbf{Notwendige Bedingungen:} Ein Graph mit Hamiltonkreis ist zusammenhängend und hat keinen Knoten mit Grad 1. Ein Graph mit Hamiltonweg ist zusammenhängend und hat höchstens zwei Knoten mit Grad 1.

\textbf{Wichtiger Unterschied Euler vs.\ Hamilton:}
\begin{itemize}[leftmargin=1.5em]
    \item Für Euler: exaktes notwendiges \emph{und} hinreichendes Kriterium bekannt (Gradparität).
    \item Für Hamilton: kein notwendiges und hinreichendes Kriterium bekannt - eines der schwierigsten ungelösten Probleme der Mathematik (eng verwandt mit dem P-NP-Problem).
\end{itemize}

\newpage
% ─────────────────────────────────────────────────────────────────────────────
\subsection{Bipartite und planare Graphen}

\definition{\textbf{Bipartiter Graph:} Ein Graph $G = (V,E)$ heißt \textbf{bipartit}, wenn sich die Knotenmenge in zwei disjunkte Teilmengen $V = U \cup W$ aufteilen lässt, sodass jede Kante genau einen Endknoten in $U$ und einen in $W$ hat. Äquivalent: $G$ enthält keinen Kreis \emph{ungerader Länge}.}

\textbf{Spezialfall Complete Bipartite Graph $K_{m,n}$:} Jeder Knoten aus $U$ ($|U|=m$) ist mit jedem Knoten aus $W$ ($|W|=n$) verbunden; $|E| = m \cdot n$.

\definition{\textbf{Planarer Graph:} Ein Graph heißt \textbf{planar}, wenn er sich in der Ebene so zeichnen lässt, dass sich keine zwei Kanten überkreuzen. (Ein Graph ist unabhängig von seiner grafischen Darstellung - Planarität ist eine Eigenschaft des Graphen, nicht der Zeichnung.)}

\textbf{Satz von Kuratowski:} Ein Graph ist genau dann planar, wenn er keinen der folgenden beiden Graphen als \textbf{Minor} enthält:
\begin{itemize}[leftmargin=1.5em]
    \item $K_5$: vollständiger Graph auf 5 Knoten (jeder Knoten mit jedem verbunden)
    \item $K_{3,3}$: vollständiger bipartiter Graph mit je 3 Knoten auf beiden Seiten
\end{itemize}

\textbf{Grundtypen von Netzwerken:}
\begin{itemize}[leftmargin=1.5em]
    \item \textbf{Complete Network:} Jeder Knoten mit jedem verbunden; die Knoten bilden eine \emph{Clique}. $|E| = \binom{|V|}{2}$.
    \item \textbf{Star Network:} Genau ein Zentralknoten verbunden mit allen anderen.
    \item \textbf{Tree Network:} Zwischen je zwei Knoten existiert genau ein Weg; kreisfrei; $|V| = |E| + 1$; alle Kanten sind Brücken.
    \item \textbf{Ring Network:} Knoten kreisförmig angeordnet, jeder mit $k$ Nachbarn verbunden. Complete Network ist Spezialfall ($k = |V|-1$).
    \item \textbf{Grid Network (Torus):} Schachbrettartige Anordnung mit horizontalen und vertikalen Verbindungen, periodisch an den Rändern.
    \item \textbf{Small-World Network:} Im Wesentlichen ein Ring Network mit zufälligen Zusatzkanten (kurze Pfade trotz hoher lokaler Vernetzung).
    \item \textbf{Power/Free-scale Network:} Wenige Knoten mit sehr hohem Grad (Hubs), viele Knoten mit sehr geringem Grad (Potenzgesetz-Gradverteilung).
\end{itemize}

\textbf{Wurzelbäume:} Ein \textbf{Wurzelbaum} ist ein Tree Network mit einem ausgezeichneten Knoten (\emph{Wurzel}). Knoten zwischen einem Knoten $y$ und der Wurzel heißen \emph{Vorgänger} von $y$ (direkter Vorgänger: \emph{Vater}); Knoten unterhalb heißen \emph{Nachfolger} (direkter Nachfolger: \emph{Sohn}). Wurzelbäume können binär kodiert werden (1 = Abstieg, 0 = Aufstieg), und diese Binärkodierung ist hexadezimal darstellbar.

\textbf{Merke:} Bipartite Graphen sind nicht notwendigerweise planar; planare Graphen sind nicht notwendigerweise bipartit.

\newpage